%% =========================================================================
\section{Ensemble Averaging and the Reduction Landscape}
\label{sec:ensemble}
%% =========================================================================

The orbit-specificity barrier---the difficulty of transferring
population statistics to one deterministic orbit---can be circumvented
by averaging \emph{over the orbit parameter itself}.  This section
develops the ensemble framework: we parametrize the EM construction by
its starting point, average over squarefree starting points via the
Chinese Remainder Theorem, and locate the orbit hypothesis CME within
the resulting landscape.

\medskip\noindent\textbf{What is new vs.\ what is infrastructure.}\enspace
The generalized EM sequences, the CRT propagation framework, and the
reduction landscape assembled around CME~$\Rightarrow$~MC
(Theorem~\ref{thm:cme-ccsb}) are original contributions of this
formalization.  The decorrelation--variance--concentration
chain (\S\ref{sec:decorr-chain}) is proved with explicit constants
($C = 2$, Theorem~\ref{thm:variance-bound}).  The sieve density function
and Buchstab weight (\S\ref{sec:first-moment}) formalize known sieve
theory.  The spectral routes that grew out of this framework---the
Weyl hitting bridge, the generalized MC argument, and the visit
equidistribution identity---are also original; they are developed
separately in Section~\ref{sec:spectral-routes}.

Sections~\ref{sec:walk}--\ref{sec:character} worked with the standard orbit
$s = 2$.  We now exploit the starting-point parameter~$s$ introduced in
\eqref{eq:gen-em}: the \emph{ensemble} is the family of EM trajectories
$\{\Prod{n} : n \text{ squarefree}\}$.  The structural properties
established in Section~\ref{sec:walk}---divisibility, monotonicity,
injectivity, squarefreeness, and the death channel
equivalence---hold for every starting point and are used throughout.

\begin{roadmap}
\begin{enumerate}[nosep,leftmargin=*]
\item[\S\ref{sec:crt-equidist}] CRT equidistribution and propagation --- ensemble residues equidistribute and propagate across steps.
\item[\S\ref{sec:decorr-chain}] Decorrelation--variance--concentration chain --- StepDecorrelation $\Rightarrow$ character sum cancellation for a.a.\ starting points.
\item[\S\ref{sec:first-moment}] Sieve structure and first moments --- ensemble average of $1/\mathrm{genSeq}$ is positive; parity and CRT bounds.
\item[\S\ref{sec:dsl-route}] The orbit hypothesis --- CME seen from the ensemble; the retired population layer.
\item[\S\ref{sec:landscape}] The reduction landscape --- the complete dependency graph of the formalization.
\end{enumerate}
\end{roadmap}

\subsection{CRT Equidistribution and Propagation}
\label{sec:crt-equidist}

The ensemble is the collection of generalized EM sequences indexed by
squarefree starting points $n \in [1, X]$.  For each prime~$q$, we
study how the residues $\Prod{n}(k) \bmod q$ and
$\seq{n}(k) \bmod q$ distribute as $n$ ranges over
squarefree integers.

\begin{definition}[\defname{Accumulator density}]
For prime~$q$, step~$k$, and residue class $a \in (\ZZ/q\ZZ)^\times$,
define the \emph{accumulator density}:
\[
  \rho_k(a) \;=\; \lim_{X \to \infty} \frac{|\{n \in [1,X] : n
    \text{ squarefree},\; \Prod{n}(k) \equiv a \pmod{q}\}|}
    {|\{n \in [1, X] : n \text{ squarefree}\}|}.
\]
\end{definition}

\begin{remark}[\defname{SquarefreeResidueEquidist} (\abbr{SRE})]
Squarefree integers equidistribute among non-zero residue classes
modulo any prime~$q$: $\rho_0(a) = 1/(q{-}1)$ for all
$a \in (\ZZ/q\ZZ)^\times$.
This is a standard consequence of the inclusion-exclusion formula for
squarefree counting (a known result, not yet in Mathlib).
\end{remark}

SRE provides the base case of a CRT induction:

\begin{property}[\defname{CRTPropagationStep}]
For all primes~$q$ and steps~$k$: if $\rho_k$ is uniform on
$(\ZZ/q\ZZ)^\times$, then so is $\rho_{k+1}$.
\end{property}

The structural basis for CRT propagation is the \emph{CRT multiplier
invariance}
(\lean{EM/Group/CRT.lean}{crt\_multiplier\_invariance}): the generalized
multiplier $\seq{n}(k) = \minFac(\Prod{n}(k) + 1)$
depends on the full integer $\Prod{n}(k) + 1$, not just its residue
modulo~$q$.  By CRT, knowing $\Prod{n}(k) \bmod q$ does not
constrain $\Prod{n}(k) \bmod r$ for any other prime~$r$, so the
factoring outcome is ``$q$-blind.''  This decorrelates the mod-$q$
coordinate of the accumulator at step~$k$ from the multiplier at
step~$k$.

The propagation proceeds by induction on~$k$:

\begin{theorem}[\code{sre\_crt\_implies\_accum\_equidist} \superseded]
\label{thm:sre-crt}
$\mathrm{SRE} + \mathrm{CRTPropagationStep}$ implies
accumulator equidistribution:
$\rho_k(a) = 1/(q{-}1)$ for all~$k$, $q$,
$a \in (\ZZ/q\ZZ)^\times$.
\end{theorem}

\noindent This implication was machine-verified, but the route is dead:
both hypotheses are RED --- CRTPropagationStep preserves the false AEP,
and AEP itself fails by absorption (Dead End~\#137).  The theorem was
deleted in the Session-288 consolidation (see git history and
\code{EM/Archive/Ensemble/CRTArchive.lean} (local archive), which retains the
surviving conditional statements of this chain).

From accumulator equidistribution, a population transfer bridge gives
\emph{ensemble multiplier equidistribution}~(EME): the density of
squarefree $n$ with $\seq{n}(k) \equiv a \pmod{q}$ tends
to $1/(q{-}1)$ for each non-zero~$a$.

\paragraph{CRT independence and joint equidistribution.}
A deeper structural property concerns the \emph{joint} distribution of
residues at different steps.  The \emph{marginal product} version---the
product $\rho_j(a) \cdot \rho_k(b)$---is trivially $1/(q{-}1)^2$ from
accumulator equidistribution.  The \emph{genuine joint}
version---counting squarefree~$n$ with $\Prod{n}(j) \equiv a$ and
$\Prod{n}(k) \equiv b$ \emph{simultaneously}---is not trivial: it
requires genuine independence of the mod-$q$ coordinates at different
steps, formalized as \defname{SquarefreeCRTIndependence}
(\code{EM/Ensemble/CRTFreedom.lean}).

The CRT freedom framework proves:

\begin{theorem}[\archived{EM/Archive/Ensemble/CRTFreedomArchive.lean}{scrti\_sre\_implies\_aep}]
$\mathrm{SCRTI} + \mathrm{SRE}$ implies
$\mathrm{AccumulatorEquidistPropagation}$.
\end{theorem}

\noindent This provides a clean decomposition: SRE is the base case
(a consequence of standard analytic number theory),
SquarefreeCRTIndependence is the inductive step (CRT structure), and
AccumulatorEquidistPropagation is the output.

\paragraph{Standard analytic number theory (ANT), and a correction.}
\label{sec:ant-def}
Throughout this paper, \textbf{ANT} denotes Dirichlet's theorem on
primes in arithmetic progressions together with sieve estimates for
shifted squarefree integers.  Earlier drafts claimed that ANT implies
\defname{MinFacResidueEquidist}~(\abbr{MFRE})---equidistribution of
$\minFac(m) \bmod q$ among $q$-rough shifted squarefree integers---and
hence PopulationEquidist~(PE).  \textbf{This is false} (Dead
End~\#160).  For a prime~$p$ the set $\{m : \minFac(m) = p\}$ has
density $w_p = p^{-1}\prod_{r<p}(1-1/r)$, the weights telescope
($w_p = c(p) - c(p{+}1)$ with $c(n) = \prod_{r<n}(1-1/r) \to 0$), and
the density of the class $\minFac \equiv a \pmod q$ among $q$-rough
integers is the \emph{convergent} series $\sum_{p \equiv a} w_p$
(\lean{EM/Population/HeadDomination.lean}{tendsto\_classCount\_div}).
The unshifted MFRE is therefore \emph{equivalent} to the identity
$\sum_{p\equiv a} w_p = c(q{+}1)/(q-1)$ for every class
(\lean{EM/Population/HeadDomination.lean}{roughLPFEquidist\_iff}), and
Dirichlet's theorem constrains only the tail of that series, which
carries no mass; the value is decided by the first primes above~$q$,
whose classes receive more than their share.  The shifted version and
PE fail identically with the sieve weights $g(r)=r/(r^2-1)$; so does
UniformConductorEquidist, whose conductor-$1$ clause is the unshifted
statement (\lean{EM/Adelic/UniformConductor.lean}{uce\_implies\_roughLPFEquidist}).
The whole population layer---PE, MFRE, the transfer lemma
$\mathrm{DSL} = (\mathrm{PE} \Rightarrow \mathrm{CME})$ and the chains
built on them---is archived
(\code{EM/Archive/Population/PopulationEquidistArchive.lean} (local archive)); the
record is \S\ref{sec:ant-formalization} and
\code{docs/pe\_dsl\_retirement.md}.  What ANT does deliver, now
unconditionally, is the entry of the analytic chain: Karamata's
Tauberian theorem (\code{EM/IK/Karamata.lean}) gives Mertens' theorem
in progressions in asymptotic form, and prime power stripping and
Abel summation carry it to
$\sum_{p \leq x,\, p \equiv a} 1/p \sim (\log\log x)/\varphi(q)$
(\code{EM/IK/Tauberian.lean}, \code{EM/IK/AbelChain.lean}).

\paragraph{The bag-conditioned multiplier law.}
\label{sec:bag-law}
The population that actually contains the Euclid numbers is not the
$q$-rough integers but the progression $m \equiv 1 \pmod{P}$, $P$ the
current accumulator; there the least prime factor is automatically outside
the bag, and its law is again exact.  For a prime $p \nmid P$,
\[
  \lim_{X\to\infty}
  \frac{\#\{m \le X : m \equiv 1\ (P),\ \minFac m = p\}}
       {\#\{m \le X : m \equiv 1\ (P)\}}
  \;=\;
  \frac1p \prod_{\substack{r<p,\ r\nmid P}} \Bigl(1-\frac1r\Bigr)
\]
(\lean{EM/Population/BagConditionedLaw.lean}{tendsto\_bagClass\_div\_ap}):
on the progression, $\minFac m = p$ reduces to $p \mid m$ together with
coprimality to the product $N'$ of the primes below~$p$ outside the bag
(\lean{EM/Population/BagConditionedLaw.lean}{minFac\_eq\_iff\_on\_ap}), the
two congruences combine by CRT into one class mod~$pP$, and along that
class the coprimality to $N'$ occurs exactly $\varphi(N')$ times per block
of~$N'$
(\lean{EM/ForMathlib/CoprimeAffineBlock.lean}{card\_coprime\_affine\_block}).
The law is head-dominated by the least primes outside the bag, exactly as
before, and it has one clean corollary: if $q$ is the \emph{least} prime
not dividing~$P$, the product is empty and the relative density of
$\{\minFac = q\}$ is exactly $1/q$, whatever the residue classes of $q$
and whatever the bag
(\lean{EM/Population/BagConditionedLaw.lean}{bagWeight\_least\_missing},
 \lean{EM/Population/BagConditionedLaw.lean}{tendsto\_least\_missing\_div\_ap}).
This is Shanks' heuristic with its correct, biased, bag-dependent law,
proved at the population level; it also says why fixed-modulus
equidistribution was never the relevant object.  What the orbit inherits
from it remains heuristic: $P_n+1$ is one member of the progression, not
a random one.

\subsection{The Decorrelation--Variance--Concentration Chain}
\label{sec:decorr-chain}

The ensemble character energy is the key object connecting CRT
equidistribution to MC\@.  For a character~$\chi$ and starting
point~$n$, the \emph{multiplier character energy} at length~$K$ is
\[
  E(n, K) \;=\; \Bigl\|\sum_{k < K} \chi\bigl(\seq{n}(k) \bmod q\bigr)\Bigr\|^2.
\]
This is an energy in the character spectrum: $E(n, K)$ lies between $0$
(perfect cancellation) and $K^2$ (no cancellation).  CCSB asks that
$E(n, K) = o(K^2)$ for the standard orbit $n = 2$; the ensemble
framework asks about the distribution of $E(n, K)$ over squarefree~$n$.

\begin{property}[\defname{StepDecorrelation} (\abbr{SD})]
For every prime~$q$, character~$\chi$ with $\|\chi(a)\|^2 \leq 1$, and
steps $j < k$: the ensemble cross-term vanishes:
\[
  \lim_{X \to \infty}
  \mathbb{E}_{n \leq X}^{\mathrm{sqf}}
  \bigl[\mathrm{Re}\bigl(\chi(m_j) \cdot \overline{\chi(m_k)}\bigr)\bigr]
  \;=\; 0,
\]
where $m_j = \seq{n}(j) \bmod q$ and the expectation is
over squarefree~$n$.
\end{property}

\noindent\emph{As stated, SD is false} (Dead End~\#156, witnessed by
\lean{EM/Ensemble/UncenteredRefutations.lean}{not\_stepDecorrelation}):
the constant function $\chi \equiv 1$ satisfies the side condition and
has correlation identically~$1$; the same defect kills the variance
bound, the concentration statement, the four-point and second-moment
hypotheses, and TWD below.  The intended statement is the
\emph{centered} one, for nontrivial Dirichlet characters; at fixed steps
$j,k$ that statement is itself of the head-dominated fixed-step type of
Dead End~\#157, so it is not pursued.  The chain that follows is kept
because it is correct conditional mathematics and its bridges are
reusable.

Step Decorrelation says that the character values at different steps are
asymptotically uncorrelated when averaged over starting points.  CRT
independence provides the structural basis: between steps $j$ and $k$,
each ``$+1$ shift'' in the Euclid construction decorrelates the mod-$q$
information.

\begin{property}[\defname{CharSumVarianceBound} (\abbr{CSVB})]
There exists $C > 0$ such that for all $q$, $\chi$, and $K$:
$\mathbb{E}_{n}^{\mathrm{sqf}}[E(n, K)] \leq C \cdot K$
eventually in~$X$.
\end{property}

\begin{theorem}[\lean{EM/Ensemble/PT.lean}{decorrelation\_implies\_variance\_proved}]
\label{thm:variance-bound}
$\mathrm{SD} \;\Longrightarrow\; \mathrm{CSVB}$ with $C = 2$.
\end{theorem}

\begin{proof}
By induction on~$K$.

\emph{Base case.}  $E(n, 0) = 0 \leq 0$ for all~$n$.

\emph{Inductive step.}  The energy recurrence gives
\[
  E(n, K{+}1) \;=\; E(n, K) + \|\chi(m_K)\|^2 +
  2\,\mathrm{Re}\Bigl(\sum_{j < K} \chi(m_j) \cdot \overline{\chi(m_K)}\Bigr).
\]
Taking ensemble averages:
\begin{itemize}[nosep]
\item \textbf{Diagonal}: $\mathbb{E}[\|\chi(m_K)\|^2] \leq 1$ since
  $\|\chi(a)\|^2 \leq 1$ pointwise.
\item \textbf{Inductive term}: $\mathbb{E}[E(n, K)] \leq 2K$ by the
  inductive hypothesis (for $X \geq X_0^{(K)}$).
\item \textbf{Cross-terms}: for each $j < K$, SD gives
  $|\mathbb{E}[\mathrm{Re}(\chi(m_j)\overline{\chi(m_K)})]|
  < 1/(2(K{+}1))$ for~$X$ large enough.
  Summing the $K$ terms: $2|\text{cross sum}| \leq 1$.
\end{itemize}
Combining: $\mathbb{E}[E(n, K{+}1)] \leq 2K + 1 + 1 = 2(K{+}1)$.
The threshold $X_0^{(K+1)}$ is the maximum of $X_0^{(K)}$ and the $K$
individual cross-term thresholds.
\end{proof}

The constant $C = 2$ is sharp for this method: each step contributes
at most $1$ from the diagonal and at most $1$ from the cross-terms.
For comparison, independent random variables with $|\chi(m)| = 1$ give
$\mathbb{E}[E(n, K)] = K$ exactly (perfect cancellation of
cross-terms); $C = 2$ allows the cross-terms to contribute up to their
maximum.

\begin{property}[\defname{EnsembleCharSumConcentration}]
For every $q$, $\chi$, and $\varepsilon > 0$: the density of
squarefree $n \leq X$ with $E(n, K) > (\varepsilon K)^2$ tends to~$0$
as $X \to \infty$, uniformly in~$K$ past a threshold.
\end{property}

\begin{theorem}[\lean{EM/Ensemble/Decorrelation.lean}{char\_variance\_implies\_concentration\_proved}]
CSVB implies EnsembleCharSumConcentration.
\end{theorem}

\begin{proof}[Proof sketch]
Markov's inequality.  Fix $\varepsilon > 0$.  The density of squarefree
$n$ with $E(n, K) > (\varepsilon K)^2$ is at most
$\mathbb{E}[E(n, K)] / (\varepsilon K)^2 \leq 2K / (\varepsilon^2 K^2)
= 2/(\varepsilon^2 K)$, which tends to $0$ as $K \to \infty$.
\end{proof}

\begin{theorem}[\lean{EM/Ensemble/Decorrelation.lean}{char\_concentration\_implies\_cancellation}]
EnsembleCharSumConcentration $\Longrightarrow$ for almost all
squarefree~$n$, the character energy $E(n, K) = o(K^2)$.
\end{theorem}

\begin{proof}[Proof sketch]
The set of $n$ with $E(n, K) > (\varepsilon K)^2$ \emph{for all}~$K$
is contained in the set for any specific~$K_0$.  Choosing $K_0$
large enough (from concentration) makes this density arbitrarily
small.
\end{proof}

\noindent The three theorems compose into a single proved reduction:

\begin{theorem}[\lean{EM/Ensemble/PT.lean}{sd\_implies\_cancellation}]
\label{thm:sd-cancel}
$\mathrm{StepDecorrelation}$ alone $\Longrightarrow$ character sum
cancellation for almost all squarefree starting points.
\end{theorem}

\noindent This composition uses zero open hypotheses: StepDecorrelation is
the sole input, and the three bridges (variance, concentration,
cancellation) are all machine-verified.  Since SD as stated is false,
the theorem is a valid conditional with a false antecedent; its value
is the bridges.

\subsection{Sieve Structure and Ensemble First Moments}
\label{sec:sieve-first-moment}

The concentration chain above converts decorrelation into cancellation
via a purely analytic argument.  This subsection develops the
complementary \emph{sieve} perspective: instead of character energy,
we track the reciprocal sums $\sum 1/\seq{n}(k)$ and use parity,
CRT, and death-density arguments to establish lower bounds on the
ensemble mean.  The two threads converge on the orbit hypothesis
(\S\ref{sec:dsl-route}): the concentration chain provides the
``if decorrelated then cancellation'' implication, while the
sieve thread provides concrete evidence that the first moment
grows linearly.

\paragraph{First Moment, Sieve Density, and the Buchstab Weight}
\label{sec:first-moment}

The ensemble average of the reciprocal $1/\seq{n}(k)$
connects to classical sieve theory via the \emph{sieve density
function}:
\[
  g(r) \;=\; \frac{r}{r^2 - 1},
\]
which gives the conditional probability that a random shifted
squarefree number $m$ (with $\mu^2(m{-}1) = 1$) is divisible by a
prime~$r$.  The algebraic properties of~$g$ are fully proved in
\code{EM/Reduction/ShiftedDensity.lean}:

\begin{itemize}[nosep]
\item $g(r) > 0$ for $r \geq 2$
  (\lean{EM/Reduction/ShiftedDensity.lean}{sieveDensity\_pos}).
\item $g(r) < 1$ for $r \geq 2$
  (\lean{EM/Reduction/ShiftedDensity.lean}{sieveDensity\_lt\_one}).
\item $g(r) > 1/r$
  (\lean{EM/Reduction/ShiftedDensity.lean}{sieveDensity\_gt\_inv}):
  a density boost from automatic squarefreeness.
\item $g(r) = \tfrac{1}{2}\bigl(\tfrac{1}{r-1} + \tfrac{1}{r+1}\bigr)$
  (\lean{EM/Reduction/ShiftedDensity.lean}{sieveDensity\_partial\_frac}).
\item $g$ is strictly decreasing on $[2, \infty)$
  (\lean{EM/Reduction/ShiftedDensity.lean}{sieveDensity\_strict\_anti}).
\end{itemize}

The \emph{Buchstab weight} $\kappa(p)$ is the sieve density at the
base prime~$p$:

\begin{theorem}[\lean{EM/Ensemble/FirstMoment.lean}{buchstabWeight\_two}]
$\kappa(2) = g(2) = 2/3$.
\end{theorem}

The first-moment/variance pattern gives a concentration result for
reciprocal sums:

\begin{theorem}[\archived{EM/Archive/Ensemble/FirstMomentArchive.lean}{first\_moment\_variance\_implies\_rsd}]
First Moment + Variance Bound $\Longrightarrow$ Reciprocal Sum
Divergence: $\sum_{k < K} 1/\seq{n}(k) \to \infty$ for
almost all squarefree~$n$.
\end{theorem}

\paragraph{Positive density from the first moment alone.}
The density-1 result above requires both the first moment (linear
mean growth) and a variance bound (from pairwise step decorrelation).
A simpler route, formalized in
\archived{EM/Archive/Population/ReciprocalSumArchive.lean}{lmg\_implies\_positive\_density\_rsd},
shows that linear mean growth \emph{alone} yields a positive density
of squarefree starting points with divergent reciprocal sums.

The argument is elementary.  Each term $1/\seq{n}(k) \leq 1/2$
(since $\seq{n}(k)$ is prime, hence $\geq 2$), so
$S_K(n) \leq K/2$ deterministically
(\lean{EM/Population/ReciprocalSum.lean}{recipPartialSum\_le\_half\_K}).
If the ensemble mean satisfies $\mathbb{E}[S_K] \geq \kappa K$ (the
\textbf{LinearMeanGrowth} hypothesis, which follows from
FirstMomentStep), a partition of the ensemble into
$\{S_K \geq M\}$ and $\{S_K < M\}$ gives
\[
  \mathrm{density}\{S_K \geq M\}
  \;\geq\;
  \frac{\kappa K - M}{K/2 - M}
  \;\xrightarrow{K \to \infty}\;
  2\kappa.
\]
Taking $\delta = \kappa/2$ yields a 1-hypothesis route:
\[
  \mathrm{LMG}
  \;\Longrightarrow\;
  \text{PositiveDensityRSD},
\]
compared to the 2-hypothesis route
PSD~$+$~LMG~$\Rightarrow$~AASRSD (density~1).
With $\kappa \geq 1/3$ from the Buchstab weight at $p=2$
(\lean{EM/Ensemble/FirstMoment.lean}{kappaPartial\_three}),
the positive density is at least $\delta = 1/6$.

The chain is now fully closed to a single hypothesis.
Linearity of the ensemble average
(\lean{EM/Ensemble/FirstMoment.lean}{ensembleAvg\_sum\_range})
gives:

\begin{theorem}[\archived{EM/Archive/Ensemble/FirstMomentArchive.lean}{first\_moment\_step\_implies\_positive\_density\_rsd}]
\label{thm:fms-pdrsd}
$\mathrm{FirstMomentStep}(\kappa) \;\Longrightarrow\;
\mathrm{PositiveDensityRSD}$, for any $\kappa > 0$.
\end{theorem}

\noindent The proof composes two links, each machine-checked before
archival: $\mathrm{FirstMomentStep} \Rightarrow \mathrm{LMG}$ (via
$\mathbb{E}[\sum] = \sum\mathbb{E}$ and a finite-max threshold
argument) and $\mathrm{LMG} \Rightarrow \mathrm{PositiveDensityRSD}$
(the partition argument above); the full three-link chain is witnessed
by \archived{EM/Archive/Ensemble/FirstMomentArchive.lean}{weak\_mc\_landscape}.
The route is conditional on RED hypotheses, however: FirstMomentStep is
RED~\#9 and LMG is RED~\#8 --- Session~201 computations indicate
$\mathbb{E}[S_K] \approx C\log\log K$, sublinear rather than linear ---
so these are valid implications whose antecedents are likely false,
retained in the archive.

\paragraph{Parity structure and unconditional bounds.}
The mod-2 structure of the generalized EM sequence yields
unconditional results that partly close the weak-MC chain.

\begin{theorem}[\lean{EM/Ensemble/FirstMoment.lean}{genSeq\_ge\_three}]
\label{thm:genSeq-ge-three}
For $k \geq 1$ and $n \geq 1$, the generalized EM prime satisfies
$\seq{n}(k) \geq 3$.
\end{theorem}

\begin{proof}
The product $\Prod{n}(1) = n \cdot \minFac(n{+}1)$ is always even:
if $n$ is even the product is even; if $n$ is odd then $n{+}1$ is
even, so $\minFac(n{+}1) = 2$
(\lean{EM/Ensemble/FirstMoment.lean}{genProd\_one\_even}).
For $k \geq 1$, $\Prod{n}(1) \mid \Prod{n}(k)$ by
\code{genProd\_dvd\_genProd}, so $\Prod{n}(k)$ is even and
$\Prod{n}(k){+}1$ is odd
(\lean{EM/Ensemble/FirstMoment.lean}{genProd\_succ\_odd}).
Hence $\minFac(\Prod{n}(k){+}1) \neq 2$
(by \code{Nat.minFac\_eq\_two\_iff}), and since
$\seq{n}(k)$ is prime with $\seq{n}(k) \geq 2$, we get
$\seq{n}(k) \geq 3$.
\end{proof}

This tightens the deterministic bound from $S_K(n) \leq K/2$
to $S_K(n) \leq 1/2 + (K{-}1)/3$ for $K \geq 1$.

\begin{theorem}[\lean{EM/Ensemble/FirstMoment.lean}{ensembleAvg\_k0\_ge\_quarter}]
\label{thm:k0-quarter}
For $X$ with $\mathrm{sqfreeCount}(X) > 0$:
$\mathbb{E}[1/\seq{\cdot}(0)] \geq 1/4$.
\end{theorem}

\begin{proof}
For odd squarefree $n$, $n{+}1$ is even, so
$\seq{n}(0) = \minFac(n{+}1) = 2$ and $1/\seq{n}(0) = 1/2$.
The density of odd squarefree integers among all squarefree
integers is $\geq 1/2$: the halving map $m \mapsto m/2$ injects
even squarefree into odd squarefree
(\lean{EM/Ensemble/FirstMoment.lean}{odd\_sf\_card\_ge\_half}),
giving $\#\{\text{even sf}\} \leq \#\{\text{odd sf}\}$.
So $\mathbb{E}[1/\seq{\cdot}(0)] \geq (1/2) \cdot (1/2) = 1/4$.
\end{proof}

\noindent This is \emph{unconditional}---no open hypothesis is used.

A weaker-than-FirstMomentStep hypothesis suffices for the full chain:
\textbf{StepMeanLowerBound}$(c)$
(\lean{EM/Ensemble/FirstMoment.lean}{StepMeanLowerBound})
asserts only that $\mathbb{E}[1/\seq{\cdot}(k)] \geq c$ eventually
for each~$k$, without requiring convergence to a specific limit.
The $k = 0$ case is proved unconditionally with $c = 1/4$
(\lean{EM/Ensemble/FirstMoment.lean}{smlb\_k0\_unconditional}).
The chain $\mathrm{SMLB} \Rightarrow \mathrm{LMG} \Rightarrow
\mathrm{PositiveDensityRSD}$ is fully machine-verified
(\archived{EM/Archive/Ensemble/FirstMomentArchive.lean}{smlb\_implies\_positive\_density\_rsd}).
If FirstMomentStep$(\kappa)$ holds, then $\kappa \geq 1/4$
(\archived{EM/Archive/Ensemble/FirstMomentArchive.lean}{fms\_kappa\_ge\_quarter},
via \code{ge\_of\_tendsto}).

For $k = 1$, the CRT structure gives further traction: for squarefree
$n \equiv 1 \bmod 6$, the sequence $\seq{n}(1) = 3$
(\lean{EM/Ensemble/FirstMoment.lean}{genSeq\_one\_of\_mod6}),
contributing $1/3$ to the sum.  The injection argument extends to
factor~3: $\#\{\text{sf}\} \leq 4 \cdot \#\{\text{odd sf coprime to 3}\}$
(\lean{EM/Ensemble/FirstMoment.lean}{sf\_le\_four\_coprime6}).
The density hypothesis \textbf{Mod6DensityLB}---that the fraction
of $1$-mod-$6$ squarefree among all squarefree exceeds $1/8$---is
proved unconditionally
(\lean{EM/Ensemble/FirstMoment.lean}{mod6\_density\_lb\_proved})
via a sieve argument: non-squarefree $n \equiv 1 \bmod 6$ must have
some $d^2 \mid n$ with $d \geq 5$ coprime to~$6$; since
$d^2 \equiv 1 \bmod 6$ (CRT), the per-$d$ count is $\leq X/(6d^2)+1$
and the telescoping bound $\sum_{d \geq 5} 1/d^2 \leq 1/4$ gives
the total non-squarefree count $\leq X/24 + \sqrt{X}$.
Combined with the injection-based lower bound
$\#\{n \equiv 1 \bmod 6\} \geq X/6$, inclusion-exclusion on
multiples of~$4$ and~$9$ for the squarefree count yields the result
for $X \geq 10{,}000$.
SMLB at $k = 1$ then follows unconditionally with $c = 1/24$
(\lean{EM/Ensemble/FirstMoment.lean}{mod6\_density\_implies\_smlb\_k1}).

The CRT analysis extends to $k = 2$: for $n \equiv 1 \bmod 6$, we have
$\Prod{n}(2) = 6n$
(\lean{EM/Ensemble/FirstMoment.lean}{genProd\_two\_of\_mod6})
and $\seq{n}(2) = \minFac(6n{+}1)$.  When additionally
$n \equiv 4 \bmod 5$ (i.e., $n \equiv 19 \bmod 30$ by CRT), we get
$5 \mid (6n{+}1)$ while $2 \nmid (6n{+}1)$ and $3 \nmid (6n{+}1)$,
so $\seq{n}(2) = 5$
(\lean{EM/Ensemble/FirstMoment.lean}{genSeq\_two\_of\_mod30}).
Thus for $n \equiv 19 \bmod 30$ squarefree, the first three
generalized EM primes are $2, 3, 5$ deterministically.

A \textbf{PartialSMLB}$(c, K_0)$ hypothesis---requiring
$\mathbb{E}[1/\seq{\cdot}(k)] \geq c$ only for $k \leq K_0$---yields a
fixed mean lower bound $\mathbb{E}[S_{K_0+1}] \geq c(K_0+1)$ via
linearity
(\lean{EM/Ensemble/FirstMoment.lean}{partial\_smlb\_implies\_mean\_lower\_bound}).
PartialSMLB$(1/4, 0)$ holds unconditionally
(\lean{EM/Ensemble/FirstMoment.lean}{partial\_smlb\_zero\_unconditional}),
giving a concrete (though non-divergent) threshold at one step.

\paragraph{The death density bridge.}
\label{sec:death-density}
The mod-6 analysis above proves SMLB at~$k=0$ and~$k=1$
unconditionally, but extending to all~$k$ requires a structural
argument.  A \emph{death density bridge} generalizes the single-prime
approach: when $\Prod{n}(k) \equiv -1 \pmod{q}$ for a prime~$q$, the
next generalized multiplier satisfies $q \mid \Prod{n}(k)+1$, so
$\seq{n}(k) = \minFac(\Prod{n}(k)+1) \leq q$
(\lean{EM/Ensemble/CRT.lean}{genSeq\_le\_of\_genProd\_neg\_one}).
Combined with $\seq{n}(k) \geq 3$ for $k \geq 1$
(Theorem~\ref{thm:genSeq-ge-three}), this gives $1/\seq{n}(k)
\geq 1/q$ on the ``death fiber''
$\{n : \Prod{n}(k) \equiv -1 \pmod{q}\}$.

\begin{theorem}[\lean{EM/Ensemble/CRT.lean}{ensembleAvg\_ge\_death\_density}]
\label{thm:death-density-bridge}
For any prime $q \geq 3$, step $k \geq 1$:
\[
  \mathbb{E}\bigl[1/\seq{\cdot}(k)\bigr]
  \;\geq\;
  \frac{\mathrm{density}\{\Prod{n}(k) \equiv -1 \!\!\pmod{q}\}}{q}.
\]
\end{theorem}

\begin{proof}
Split the ensemble sum by the death fiber.  On the fiber,
each term $\geq 1/q$; on the complement, each term $\geq 0$.
The sum on the fiber is therefore at least
$\#\{\text{fiber}\}/q$.  Dividing by the total squarefree
count gives the density bound.
\end{proof}

\noindent At $q = 3$: since $-1 \equiv 2 \pmod{3}$ and $3 \mid
\Prod{n}(k)+1$ forces $\seq{n}(k) = 3$ (odd
and divisible by~$3$), the bound becomes
$\mathbb{E}[1/\seq{\cdot}(k)] \geq d_2/3$ where~$d_2$ is the
density of genProd~$\equiv 2 \pmod{3}$
(\lean{EM/Ensemble/CRT.lean}{ensembleAvg\_ge\_mod3\_density}).
The general bridge unifies with the mod-3 case:
$\mathrm{AccumMod3LB}(c) \Leftrightarrow
\mathrm{DeathDensityLB}(3, c)$
(\code{accumMod3LB\_iff\_deathDensity3}, deleted in the Session-288
consolidation --- AccumMod3LB is RED~\#6, false since the death density
decays to zero; the surviving statements are in
\code{EM/Archive/Ensemble/CRTArchive.lean} (local archive)).

A positive lower bound on the death density at any single
prime~$q \geq 3$ for all steps~$k$ yields the full chain:

\begin{theorem}[\archived{EM/Archive/Ensemble/CRTArchive.lean}{death\_density\_implies\_prsd}]
\label{thm:death-density-prsd}
For any prime $q \geq 3$ and constant $c > 0$:
\[
  \mathrm{DeathDensityLB}(q, c)
  \;\Longrightarrow\;
  \mathrm{SMLB}\bigl(\min(1/4,\, c/q)\bigr)
  \;\Longrightarrow\;
  \mathrm{PositiveDensityRSD}.
\]
\end{theorem}

\noindent The four routes are assembled in
\archived{EM/Archive/Ensemble/CRTArchive.lean}{ewe\_landscape\_extended}.

\paragraph{The absorption mechanism.}
\label{sec:absorption}
While the death density bridge (Theorem~\ref{thm:death-density-bridge})
converts positive death density into SMLB, the death density itself
decays over time due to an \emph{absorption mechanism}.  The mod-$q$
dynamics of $\Prod{n}(k)$ has three structural properties, all
machine-verified
(\lean{EM/Ensemble/CRT.lean}{absorption\_kills\_death\_density}):

\begin{enumerate}[nosep]
\item \textbf{State~$0$ is absorbing}
  (\lean{EM/Ensemble/CRT.lean}{genProd\_mod\_zero\_absorbing}):
  if $q \mid \Prod{n}(k)$, then $q \mid \Prod{n}(k{+}j)$ for all $j \geq 0$,
  since divisibility is preserved under the recurrence
  $\Prod{n}(k{+}1) = \Prod{n}(k) \cdot \seq{n}(k)$.
\item \textbf{Death feeds absorption}
  (\lean{EM/Ensemble/CRT.lean}{death\_implies\_absorption}):
  when $\Prod{n}(k) \equiv -1 \pmod{q}$ and $\seq{n}(k) = q$,
  we get $\Prod{n}(k{+}1) = \Prod{n}(k) \cdot q \equiv 0 \pmod{q}$,
  so the orbit is absorbed.
  For $q = 3$ this always fires
  (\lean{EM/Ensemble/CRT.lean}{mod3\_death\_implies\_absorption}).
\item \textbf{Absorbed $\neq$ death}
  (\lean{EM/Ensemble/CRT.lean}{absorbed\_not\_in\_death\_class}):
  for $q \geq 2$, the absorbed state $0$ differs from the death
  state $-1$ in $\ZZ/q\ZZ$, so absorbed orbits \emph{never return} to
  the death class
  (\lean{EM/Ensemble/CRT.lean}{death\_then\_never\_death\_again}).
\end{enumerate}

\noindent Together, these show that each visit to the death class
permanently removes an orbit from future death events.
Consequently, $\mathrm{DeathDensityLB}(q, c)$ for any fixed $c > 0$
is heuristically false: the death density decays roughly as
$O\bigl((1 - 1/(q{-}1))^k\bigr)$ as orbits accumulate in the
absorbing state.  More critically, the absorption mechanism shows
that AccumulatorEquidistPropagation~(AEP) itself is false:
the correct squarefree residue density per coprime class is
$r/(r^2{-}1)$, not $1/(r{-}1)$, and even this corrected limit
fails for $k \geq 1$ because nonzero-class mass drains monotonically
to the absorbing class~$0$ (Dead Ends~\#137--\#138).
The backward dynamics chain ETA+DCTA+SRE~$\Rightarrow$~PRSD
remains formally valid but is vacuously true since
CRTPropagationStep (which preserves the AEP limit) is false.
The SMLB and DeathDensityLB routes are not affected---they
use independent, proved lower bounds.

\paragraph{The divergence hierarchy.}
\label{sec:divergence-hierarchy}
When the step mean lower bound SMLB cannot be established with a
\emph{uniform} constant~$c$ across all steps, a weaker condition
may still suffice.  Two relaxations are formalized in
\texttt{FirstMoment.lean}:

\begin{definition}[\defname{DecayingSMLB}
  (\lean{EM/Ensemble/FirstMoment.lean}{DecayingSMLB})]
A sequence $f(k) > 0$ with $\sum_k f(k) = \infty$ is a
\emph{decaying step mean lower bound} if
$\mathbb{E}[1/\seq{\cdot}(k)] \geq f(k)$ for all~$k$
(eventually in $X$).
\end{definition}

\begin{definition}[\defname{FirstMomentDivergence}
  (\lean{EM/Ensemble/FirstMoment.lean}{FirstMomentDivergence})]
$\mathbb{E}\bigl[\sum_{k<K} 1/\seq{\cdot}(k)\bigr] \to \infty$
as $K \to \infty$.
\end{definition}

\noindent The hierarchy is:
\[
  \mathrm{SMLB}(c)
  \;\Longrightarrow\;
  \mathrm{DecayingSMLB}(\text{const})
  \;\Longrightarrow\;
  \mathrm{FMD},
\]
both implications proved
(\archived{EM/Archive/Ensemble/FirstMomentArchive.lean}{smlb\_implies\_decaying\_smlb},
\lean{EM/Ensemble/FirstMoment.lean}{decaying\_smlb\_implies\_fmd}).
However, the reverse implications fail:
DecayingSMLB captures a scenario where each step mean decays but
the cumulative contribution diverges (e.g.\ $f(k) = 1/k$),
while FMD does \emph{not} imply PositiveDensityRSD---it says only
that the expected partial sum grows, not that a positive proportion
of starting points have large partial sums.  The critical gap is
exactly the passage from first-moment divergence to concentration,
which requires variance bounds
(\S\ref{sec:variance-route}).

\paragraph{Conditional MFRE (retired).}
An earlier draft refined the first-moment argument through a
\textbf{MFREConditional} hypothesis---among squarefree $m \equiv c
\pmod q$, the conditional distribution of $\minFac(m{+}1) \bmod q$ is
$1/(q{-}1) + O(1/q^2)$ uniformly in~$q$---and an
\textbf{EnsembleSelectionLemma} transferring it to the ensemble.
MFREConditional is false: for odd~$m$, $\minFac(m{+}1) = 2$, so the
class of $2 \bmod q$ carries at least two thirds of the conditional
mass, incompatible with a uniform $O(1/q^2)$ correction (Dead
End~\#160; the pair is archived in
\code{EM/Archive/Ensemble/CRTArchive.lean} (local archive)).  What survives is the
population-level CRT statement
\textbf{MinFacSelectionIndependence}
(\lean{EM/Ensemble/CRT.lean}{MinFacSelectionIndependence}): for
$q \geq 3$ and non-death classes $d_1, d_2 \neq -1$, the conditional
density of $\minFac(m{+}1) \equiv a \pmod{q}$ given
$m \equiv d_i \pmod{q}$ is asymptotically the same for $d_1$ and~$d_2$
(the classes need not be equidistributed---they are not---only
independent of the conditioning).

The CRT independence thus operates at three levels, documented in
\archived{EM/Archive/Ensemble/CRTArchive.lean}{crt\_blindness\_landscape}:
\begin{enumerate}[nosep]
\item \textbf{Pointwise} (\code{crt\_multiplier\_invariance}, proved):
  deterministic---if $P \equiv P' \pmod{r}$ for all $r \neq q$, then
  $\minFac(P{+}1) = \minFac(P'{+}1)$.
\item \textbf{Population} (MSI, open): statistical---conditional
  density of $\minFac$ mod~$q$ is independent of the mod-$q$ class.
\item \textbf{Orbit} ($=$ CME, open): along the specific EM orbit,
  the multiplier equidistributes.
\end{enumerate}

\subsection{The Orbit Hypothesis: CME in the Ensemble Picture}
\label{sec:dsl-route}

The ensemble results above are \emph{population} statements: they
concern almost all squarefree starting points.  The conjecture concerns
one orbit, $n = 2$.  Earlier drafts closed this gap with a transfer
lemma, $\mathrm{DSL} = (\mathrm{PE} \Rightarrow \mathrm{CME})$, taking as
input a population equidistribution PE believed to follow from ANT\@.
PE is false (\S\ref{sec:ant-def}, Dead End~\#160), so that lemma is
vacuous and is retired; the orbit hypothesis stands on its own as
CME (Property~\ref{def:cme}), with the machine-checked chain
CME~$\Rightarrow$~CCSB~$\Rightarrow$~MC
(\lean{EM/CME/Reduction.lean}{cme\_implies\_mc}).  What the ensemble
contributes is not an input to CME but evidence for it, and a precise
statement of the remaining gap: the population concentration
property (Theorem~\ref{thm:sd-cancel}) shows that \emph{almost all}
squarefree starting points have cancelling character energy, and
``almost all'' does not cover $n = 2$.

\begin{remark}[Why CME should hold]
\label{rem:why-dsl}
CME asserts that along the EM orbit the multiplier residues are
equidistributed on each fiber of the walk position.  Three structural
features of the EM sequence support this expectation.

\smallskip
\noindent\textbf{(1) The information bottleneck.}\enspace
At step~$n$, $\Prod{2}(n)$ has ${\sim}2^n$ bits, but the walk
position $\walkZ{2}{q}{n} = \Prod{2}(n) \bmod q$ carries only
$O(\log q)$ bits.  The smallest prime factor of $\Prod{2}(n)+1$
depends on the full ${\sim}2^n$-bit value, not on its $O(\log q)$-bit
residue.  Conditioning on $\walkZ{2}{q}{n} = c$ constrains almost
nothing about the factoring behavior of $\Prod{2}(n)+1$: the
conditional and unconditional distributions of
$\minFac(\Prod{2}(n)+1) \bmod q$ should agree, which is exactly what
CME says.  This is the same mechanism that makes iterated hash chains
pass statistical tests for independence: the forward map is easy, but
projecting an exponentially large integer onto a small residue
destroys all correlation.

\smallskip
\noindent\textbf{(2) No arithmetic obstruction.}\enspace
The products $\Prod{2}(n)$ are squarefree
(\lean{EM/Population/WeakErgodicity.lean}{prod\_squarefree}), pairwise
multiplicatively entangled (each divides the next), and each
$\Prod{2}(n)+1$ has a distinct smallest prime factor
($\seq{2}$ is injective).  There is no known
arithmetic mechanism---congruence obstruction, algebraic identity, or
Galois constraint---that would systematically bias
$\minFac(\Prod{2}(n)+1)$ toward or away from particular residue
classes mod~$q$, conditional on the walk position.  Every structural
feature that \emph{could} create such a bias (coprime cascade, CRT
propagation, cofactor identity) has been exhaustively checked and
found insufficient (\S\ref{sec:hard}).

\smallskip
\noindent\textbf{(3) The ensemble picture, honestly.}\enspace
The ensemble framework parametrizes the EM construction by squarefree
starting points.  Earlier drafts claimed here that CME holds for almost
all starting points via a proved ``StepDecorrelation~$\Rightarrow$
cancellation'' chain; that chain is a valid conditional whose
hypothesis is \emph{false} (Dead End~\#156, \S\ref{sec:decorr-chain}),
so \textbf{no almost-all statement about CME survives the audit}.  What
the ensemble does supply is unconditional structure (CRT blindness of
the multiplier, the tail identity, the bag-conditioned law of
\S\ref{sec:bag-law}) and the reminder that $2$ is one squarefree seed
among many.  This is a heuristic, not an argument; and CME being
equivalent to MC (\S\ref{sec:intro}), ``evidence for CME'' is evidence
for MC.

\smallskip
CME is therefore a \emph{typicality assertion}: the EM orbit behaves
like a generic member of its ensemble.  This is analogous to
Hooley's conditional proof of Artin's conjecture~\cite{Hooley1967},
which shows the conjecture for all integers conditional on GRH---the
hypothesis being that the specific integers of interest are not
exceptional among all integers.  The difficulty of proving CME is
the difficulty of proving that a specific deterministic sequence
satisfies a statistical law.
\end{remark}

\begin{remark}[\lean{EM/Transfer/CRTPointwise.lean}{oce\_eq\_cme}]
An alternative perspective on CME is the \emph{CRT pointwise transfer}:
condition on the walk position $c = \walkZ{2}{q}{n}$ and ask whether
the multiplier residues in each fiber are equidistributed.
This \emph{orbit-conditional equidistribution}~(OCE) turns out to be
definitionally equal to CME---proved by \code{rfl}, since the walk
position \emph{is} the residue class that CME conditions on.
Any route to OCE automatically gives CME\@.  (The associated
``CRT pointwise transfer bridge'' from a population-conditional
statement PCE to OCE is CME itself, since PCE was recorded only as a
placeholder proposition:
\lean{EM/Transfer/CRTPointwise.lean}{crtPointwiseTransferBridge\_iff\_cme}.
Likewise the ``Substitution Principle'' below is CME by definition,
\code{sp\_eq\_cme}.  Neither is a separate open point.)
\end{remark}

\begin{remark}[The tail identity and CME]\label{rem:tail-dsl}
The tail identification (Theorem~\ref{thm:tail-id}) gives a precise
structural link between the standard orbit and the ensemble.  Since
$\Prod{2}(M)$ is squarefree for all~$M$
(\lean{EM/Population/WeakErgodicity.lean}{prod\_squarefree}), the $M$-th
accumulator is itself a valid ensemble member.  The tail identity says
the orbit starting from $\Prod{2}(M)$ \emph{is} the suffix of the
standard orbit from step~$M$ onward:
\[
  \Prod{\Prod{2}(M)}(k) \;=\; \Prod{2}(M + k),
  \qquad
  \seq{\Prod{2}(M)}(k) \;=\; \seq{2}(M + k).
\]
Consequently, the character energy of the ensemble member
$n = \Prod{2}(M)$ over a window of length~$K$ equals the character
energy of the standard orbit over~$[M, M{+}K)$:
\[
  E\bigl(\Prod{2}(M),\, K\bigr)
  \;=\; E_{\mathrm{tail}}(2,\, M,\, K).
\]
The ensemble concentration result
(SD~$\Rightarrow$~CSVB~$\Rightarrow$~concentration,
Theorem~\ref{thm:sd-cancel}) says: for almost all squarefree~$n$,
$E(n, K) = o(K^2)$.  So the question ``does the standard orbit have
cancelling character energy on long windows?''\ is the same as ``does
the ensemble member $\Prod{2}(M)$ fall outside the bad set?''---for
every~$M$.

\medskip
\noindent\textbf{The Borel--Cantelli program.}\enspace
Let $B_\varepsilon(K)$ denote the set of squarefree~$n$ with
$E(n,K) > (\varepsilon K)^2$.  The first-moment (Markov) bound gives
bad density~$O(1/K)$, which is not summable over windows and
therefore insufficient for a Borel--Cantelli argument.  The
\emph{second-moment} (Chebyshev) approach resolves this: if the
population fourth moment satisfies
$\mathbb{E}_{\mathrm{sqf}}[E(n,K)^2] = O(K^2)$, then Chebyshev
gives
\[
  \Pr_{\mathrm{sqf}}\bigl[n \in B_\varepsilon(K)\bigr]
  \;\leq\; \frac{D}{\varepsilon^4 K^2},
\]
which \emph{is} summable.  Applying this to dyadic windows
$[2^i, 2^{i+1})$ with $K = 2^i$:
$\sum_i D/(\varepsilon^4 \cdot 4^i) < \infty$, so by
Borel--Cantelli, for all large~$i$ the character energy on the window
$[2^i, 2^{i+1})$ is at most $(\varepsilon \cdot 2^i)^2$.  Summing
over windows covering~$[0,N)$:
\[
  |F(c,\,[0,N))| \;\leq\;
  \sum_{i:\, 2^i < N} \varepsilon \cdot 2^i
  \;\leq\; 2\varepsilon N,
\]
which gives $|F(c,N)| = o(N)$, i.e.\ CME.

\medskip
\noindent\textbf{The sole new content: four-point PCV} (false as stated at $\chi\equiv 1$,
\code{not\_fourPointPCV}; the paragraph is kept as the record of the
intended argument).\enspace
The fourth moment $\mathbb{E}[E(n,K)^2]$ decomposes into a diagonal
contribution~$K^2$, a cross-term contribution that vanishes by PCV,
and a \emph{four-point correlation}
\[
  \Gamma_4(j_1,k_1,j_2,k_2)
  \;=\; \mathbb{E}_{\mathrm{sqf}}\bigl[
    \mathrm{Re}(\chi(m_{j_1})\overline{\chi(m_{k_1})})
    \cdot
    \mathrm{Re}(\chi(m_{j_2})\overline{\chi(m_{k_2})})
  \bigr].
\]
For distinct pairs $(j_1,k_1) \neq (j_2,k_2)$ with all four indices
distinct, CRT independence should force $\Gamma_4 \to 0$ via
factorization into two-point correlations.  If this \emph{four-point
PCV} holds, $\mathbb{E}[E(n,K)^2] = O(K^2)$, and the
Borel--Cantelli program would give CME\@.

\medskip
\noindent\textbf{The Substitution Principle (and its limits).}\enspace
An equivalent reformulation of CME, proved definitionally equal
(\code{sp\_eq\_cme}), works at the level of \emph{fiber return visits}:
the steps~$k$ where $\walkZ(q,k) = a$.  At such steps the multiplier
primes $\seq(k{+}1)$ are \emph{distinct} (\code{return\_visit\_mult\_injective})
and \emph{pairwise coprime} (\code{return\_visit\_mult\_coprime}).
The structural package is recorded in a 4-tuple witness
\code{sp\_structural\_summary}: injectivity, coprimality, primality,
and divisibility.  Four independent routes from SP to MC are assembled
in \code{all\_routes\_to\_mc\_with\_sp}.

However, the coprime cascade structure alone does \emph{not} force
character cancellation: for any $q \geq 3$, one can construct
pairwise coprime integers $N_i = p_i r_i$ in a fixed residue class
$\bmod q$ whose least prime factors all lie in the same residue class
(e.g.\ $p_i \equiv 2$, $r_i \equiv 3 \pmod{5}$ gives
$\sum \chi(\mathrm{minFac}(N_i)) = I \cdot \chi(2) = \Theta(I)$).
Thus the cancellation in CME, if true, must rely on the specific
feedback dynamics of the EM sequence---the fact that each
$\seq(n{+}1) = \mathrm{minFac}(\Prod{}(n){+}1)$ depends on all prior
terms---not merely on coprimality.  The ``Substitution Principle'' is
therefore just CME under a structurally suggestive name; the name does
not isolate a new mechanism.

\medskip
\noindent\textbf{The orbit-specificity gap.}\enspace
Even the Chebyshev bound is a statement about the \emph{density} of
bad starting points, not a pointwise statement about the specific
values~$\Prod{2}(M)$.  A density-zero set can still contain a
specific super-exponentially growing sequence.  However, the
$O(1/K^2)$ rate makes the gap much narrower: combined with CRT
typicality of $\Prod{2}(M) \bmod q$ (the residues form a walk on
$(\mathbb{Z}/q\mathbb{Z})^\times$ visiting every class cofinally, by
the proved \code{SubgroupEscape} and
\code{prime\_residue\_escape}), the gap reduces to a quantitative
equidistribution-in-residue-classes statement.  This is the sharpest
form of the orbit-specificity barrier.
\end{remark}


A structurally different family of routes---the Weyl hitting bridge,
visit equidistribution, the Lyapunov drift analysis, and tail window
decorrelation---bypasses CME entirely by working in the character
spectrum of the walk itself.  These are developed in
Section~\ref{sec:spectral-routes}.

\subsection{The Reduction Landscape}
\label{sec:landscape}

Figure~\ref{fig:reduction-graph} shows the dependency graph of the
formalization.  Every solid arrow is a machine-verified implication
with zero \code{sorry}.  Red boxes are the open hypotheses; the sole
one on the central path is CME\@.

\begin{figure}[!t]
\centering
\begin{tikzpicture}[scale=0.88, every node/.append style={transform shape},
  node distance=0.85cm and 1.3cm,
  every node/.style={font=\small},
  proved/.style={-{Stealth[length=5pt]}, thick, color=leanink},
  open/.style={-{Stealth[length=5pt]}, thick, dashed, red!70!black},
  equiv/.style={{Stealth[length=5pt]}-{Stealth[length=5pt]}, thick, color=leanink},
  box/.style={draw, rounded corners=3pt, inner sep=3pt, minimum height=1.5em,
              font=\small},
  openbox/.style={box, draw=red!70!black, fill=red!5},
  provedbox/.style={box, draw=leanink, fill=blue!5},
  goalbox/.style={box, draw=black, fill=green!8, line width=1pt},
]
% --- Row 1: CME (the open orbit hypothesis), SD, SVE <-> VE, HOD ---
\node[openbox] (cme) {\textbf{CME}};
\node[left=1.6cm of cme, provedbox] (sd) {SD};
\node[right=1.4cm of cme, provedbox] (sve) {SVE};
\node[right=0.8cm of sve, provedbox] (ve) {VE};
\node[left=1.2cm of sd, provedbox] (hod) {HOD};

% --- Row 6: CSVB, CCSB ---
\node[below=of cme, provedbox] (ccsb) {CCSB};
\node[below=of sd, provedbox] (csvb) {CSVB};

% --- Row 7: SE, SHH, Hits, DH, Weyl ---
\node[below=1.0cm of ccsb, provedbox] (hits) {Hits $-1$};
\node[right=1.2cm of hits, provedbox] (dh) {DH};
\node[right=1.0cm of dh, provedbox] (wb) {Weyl};
\node[left=1.2cm of hits, provedbox] (shh) {SHH};
\node[left=1.0cm of shh, provedbox] (se) {SE};
\node[left=1.0cm of se, provedbox] (pre) {PRE};

% --- Row 8: MC ---
\node[below=1.0cm of hits, goalbox] (mc) {\textbf{MC}};

% --- Arrows ---
% CME -> CCSB -> Hits -> MC
\draw[proved] (cme) -- (ccsb);
\draw[proved] (ccsb) -- (hits) node[midway, right, font=\tiny] {Fourier};
\draw[proved] (hits) -- (mc);

% Spectral equivalence: CCSB -> SVE <-> VE
\draw[proved] (ccsb.east) -- ++(0.3,0) |- (sve);
\draw[equiv] (sve) -- (ve);

% Decorrelation chain: SD -> CSVB -a.a.-> CCSB
\draw[proved] (sd) -- (csvb);
\draw[proved] (csvb) -- (ccsb) node[midway, below, font=\tiny] {a.a.};

% Weyl bridge: SVE -> Weyl -> DH -> MC
\draw[proved] (sve) |- (wb);
\draw[proved] (wb) -- (dh);
\draw[proved] (dh) -- (mc);

% Bootstrap: PRE -> SE -> SHH -> MC
\draw[proved] (pre) -- (se) node[midway, above, font=\tiny] {$+\MC({<}\,p)$};
\draw[proved] (se) -- (shh) node[midway, above, font=\tiny] {free};
\draw[proved] (shh) -- (mc);

% HOD -> CCSB (proved)
\draw[proved] (hod) |- (ccsb.west);


\end{tikzpicture}
\caption{Dependency graph of the main reductions.
  Solid {\color{leanink}blue} arrows: machine-verified.
  Red box: the open hypothesis CME\@.  (SD and CSVB are drawn for the
  proved bridges; SD as stated is false, Dead End~\#156.)
  All abbreviations are defined in the glossary
  (Appendix~\ref{sec:glossary}).}
\label{fig:reduction-graph}
\end{figure}

The graph reveals the architecture: multiple independent paths converge
on MC\@.
The \emph{left path} (PRE~$\to$~SE~$\to$~SHH~$\to$~MC) is the
inductive bootstrap of Section~\ref{sec:bootstrap}, proved
unconditionally---it reduces MC to producing one hit on~$-1$ at each
prime.  The \emph{central path}
(CME~$\to$~CCSB~$\to$~hits~$\to$~MC)
produces the required hit via harmonic analysis, conditional on CME\@.  The \emph{right path}
(SVE~$\leftrightarrow$~VE, Weyl~$\to$~DH~$\to$~MC) provides an
alternative spectral route.  One additional route reaches CCSB:
HOD~$\to$~CCSB (proved, via van der Corput).
All solid arrows are machine-verified; CME is the open hypothesis.
